\chapter{Glossary}\label{app:glossary}

The preceding chapters use vocabulary from three communities. The first is the proof theory of Miller and
Nadathur's book~\cite{MillerNadathur2012}, where you find sequents, uniform proofs and eigenvariables. The second
is the literature on unification of $\lambda$-terms (patterns, raising, pruning, universe levels). The third is
miniKanren, with streams, interleaving, soft cut and reification, and in particular the engine \micro{} that we
want to extend. The same word sometimes means different things in different communities (\enquote{suspension},
\enquote{level}, \enquote{scope extrusion}, \enquote{fresh}), and the same idea sometimes has three names
(GENERIC, \lp{pi}, \texttt{all}). This appendix collects all of them in one alphabetical list. Each entry gives a
short, precise definition and tells you where the concept is treated in this document and in the book.

\section{How to read this glossary}\label{sec:glossary-howto}

\paragraph{Pointers.} Most entries end with a pointer of the form
\enquote{$\rightarrow$ \cref{ch:hh}; \book{82}}. The first part is the chapter of this document where the
concept is developed. The second part is a \emph{printed} page of the book~\cite{MillerNadathur2012}; most of
these come from the book's own index. Entries that come only from the miniKanren literature or from the
\aux{} sources have no book page; they cite the paper instead. Names of \micro{} procedures and macros are set in
\texttt{typewriter} type. We describe only what exists today. The new operators are the subject of
\cref{ch:design,ch:module}.

\paragraph{Notation.} \Cref{tab:glossary-notation} lists the symbols that occur most often in the entries.

\begin{table}[ht]
\centering\small
\caption{Notation used throughout the glossary.}\label{tab:glossary-notation}
\begin{tabularx}{\textwidth}{@{}lX@{}}
\toprule
Symbol & Meaning \\
\midrule
$\seq{\Sig}{\Prog}{G}$ & sequent: signature $\Sig$, program $\Prog$, goal $G$ \book{39} \\
$\Sig;\Prog \xrightarrow{D} A$ & focused sequent: the clause $D$ is selected to prove the atom $A$ \book{44} \\
$G$, $D$ & goal formula, program clause (definite formula) \\
$\top,\wedge,\vee,\imp,\forall,\exists$ & logical constants; concrete syntax \lp{true}, \lp{&} or \lp{,}, \lp{;}, \lp{=>} or \lp{:-}, \lp{pi}, \lp{sigma} \\
$\abs{x}{t}$ & abstraction; λProlog \lp{x\ t} \\
$\subst{s}{t}{x}$ & capture-avoiding substitution of $t$ for $x$ in $s$ \book{85} \\
$=_{\beta\eta}$, $\lambda$-conversion & equality of terms modulo $\alpha$, $\beta$, $\eta$ \book{100} \\
$\tau\to\sigma$, $\typ{o}$ & function type; type of formulas \\
$Q_1x_1\dots Q_nx_n[\,\bigwedge_j t_j=s_j\,]$ & unification problem under a mixed prefix \book{105} \\
$a\mathbin{\#}t$ & nominal freshness: the atom $a$ does not occur free in $t$ \\
$\langle\phi;t\rangle$ & static closure of Ma and Friedman: term $t$ under the scope $\phi$ \\
\texttt{°}, \texttt{ª}, \texttt{§}, \texttt{δ} & \aux{} suffixes and glyphs: goal or relation, applicative form, stream, delay \\
\bottomrule
\end{tabularx}
\end{table}

\begin{pitfall}[one word, several meanings]
Several terms below are overloaded across the sources. Keep the distinctions in mind when you read the papers.
\begin{itemize}[nosep]
\item \textbf{Suspension.} In \micro{}, a suspension is a \texttt{δ}-promise, the immature part of a stream. In
  nominal unification, $\pi\cdot X$ is a permutation suspended on a variable. In Teyjus, the suspension
  $[[t,ol,nl,e]]$ is an explicit-substitution term. In ELPI, a suspended goal is a syntactic constraint.
\item \textbf{Level.} A de Bruijn level numbers a bound variable from the root inwards. A universe level, or
  universe index, is the $\forall$-nesting depth at which a variable or constant was created. ELPI makes the two
  coincide, but they are different ideas.
\item \textbf{Scope extrusion.} The book uses this name for a classical equivalence that moves an implication over
  a disjunction \book{89}. In the $\pi$-calculus it means sending a restricted name out of its scope. In
  implementation discussions it means an eigenvariable escaping into an older variable.
\item \textbf{Fresh.} In miniKanren, \texttt{fresh} is the existential quantifier and introduces logic variables.
  In αKanren, \texttt{fresh} introduces \emph{noms}, which are names rather than variables.
\item \textbf{Cut.} The cut rule of the sequent calculus has no connection to Prolog's pruning operator \lp{!}
  \book{63}.
\item \textbf{Substitution and grafting.} Instantiating a unification variable is capture-avoiding in standard
  higher-order unification. It is capture-permitting (grafting) in the Ma--Friedman and DHKP formulations.
\end{itemize}
\end{pitfall}

\section{Terms from A to Z}\label{sec:glossary-az}

\newlist{glossarylist}{description}{1}
\setlist[glossarylist]{font=\normalfont\bfseries,leftmargin=1.2em,labelindent=0pt,itemsep=0.45em,parsep=0.2em}
\begingroup\sloppy

\subsection*{A}
\begin{glossarylist}
\item[Abstract datatype] A module that exports a type constructor and the predicates over it but hides the data
  constructors, so other code can use only the exported operations (book Ch.~6). $\rightarrow$ \book{165}.
\item[Abstract logic programming language] A logic in which searching for uniform proofs is a complete proof
  procedure. This is the book's criterion for a logic to count as a logic programming language. fohc, fohh, hohc
  and hohh all qualify~\cite{MillerNadathurPfenningScedrov1991}. $\rightarrow$ \cref{ch:horn}; \book{73}.
\item[Abstraction] The term $\abs{x}{t}$ (λProlog \lp{x\ t}), the only binder of the term language. $x$ is the
  bound variable and $t$ the body. The quantifiers are the constants \lp{pi} and \lp{sigma} applied to
  abstractions. $\rightarrow$ \cref{ch:lambda}; \book{97}.
\item[Accumulation (module, signature)] \texttt{accumulate} inlines the code of other modules after renaming
  their private names. \texttt{accum\_sig} inserts the declarations of other signatures, either to re-export them
  or to declare a parameter that a later module supplies. $\rightarrow$ \book{154}, \book{155}.
\item[Administrative redex] A $\beta$-redex that a program transformation (here the CPS transformation) introduces
  itself, as opposed to one present in the source. The two-phase CPS of book Ch.~10 marks these redexes and
  reduces them afterwards. $\rightarrow$ \cref{ch:apps}; \book{257}.
\item[$\alpha$-conversion] Renaming bound variables. Terms that differ only by $\alpha$-conversion are identified,
  so the names of bound variables are \enquote{a fiction} and an implementation may not represent them at all.
  $\rightarrow$ \cref{ch:lambda,ch:lambdatree}; \book{100}, \book{183}, \book{205}.
\item[αKanren] Byrd and Friedman's nominal miniKanren~\cite{ByrdFriedman2007,Byrd2009}. \texttt{fresh}
  introduces noms, \texttt{(tie a t)} binds the nom $a$ in $t$, and \texttt{(hash a t)} asserts the freshness
  $a\mathbin{\#}t$. Unification is nominal unification. The implementation does not support disequality
  constraints. $\rightarrow$ \cref{ch:lkanren}.
\item[AND] The goal-reduction rule for $G_1\wedge G_2$: prove both conjuncts with the same signature and program.
  $\rightarrow$ \cref{ch:horn}; \book{39}.
\item[Anonymous variable] \lp{_}: each occurrence is a distinct variable, and its binding is not reported.
  $\rightarrow$ \cref{ch:horn}; \book{47}.
\item[Answer substitution] The bindings that a successful search reports for the free variables of a query. These
  variables are implicitly existential. The miniKanren counterpart is the reified answer.
  $\rightarrow$ \cref{ch:horn}; \book{43}.
\item[Applicative form] An \aux{} naming convention. The suffix \texttt{ª} marks the procedure behind a goal macro:
  \texttt{freshª} behind \texttt{fresh°}, and \texttt{andª}, \texttt{orª} behind \texttt{and°}, \texttt{or°}.
  $\rightarrow$ \cref{ch:micro}.
\item[\texttt{assume}, \texttt{all}] The λKanren goals for $D\imp G$ and $\forall x\,G$. See \emph{AUGMENT},
  \emph{GENERIC} and \emph{λKanren}. $\rightarrow$ \cref{ch:lkanren}.
\item[Atom (nominal), nom, name] In nominal logic and αKanren, a constant-like name that stands for an
  object-level variable. A nom unifies only with itself or with an unbound logic variable.
  $\rightarrow$ \cref{ch:lkanren}.
\item[Atomic formula] A term of type $\typ{o}$ whose head is a predicate symbol, that is, a constant with target
  type $\typ{o}$. $\rightarrow$ \cref{ch:horn}; \book{34}.
\item[AUGMENT] The goal-reduction rule for $D\imp G$: add $D$ to the program and prove $G$. This rule is the basis
  of hypothetical reasoning; the clause is visible only during the proof of $G$. λProlog writes \lp{D => G};
  λKanren calls it \texttt{assume}.
  \[\inferrule{\seq{\Sig}{\Prog\cup\{D\}}{G}}{\seq{\Sig}{\Prog}{D\imp G}}\;\rulename{augment}\]
  $\rightarrow$ \cref{ch:hh}; \book{39}, \book{77}.
\item[Available level] (Ma--Friedman) A de Bruijn level of a scope $\phi$ whose binder name is not rebound further
  in. A shadowed binder's level is unavailable, so $\lambda a.\lambda a.X \stackrel{?}{=} \lambda c.\lambda d.c$
  fails~\cite{MaFriedman2021}. $\rightarrow$ \cref{ch:lkanren}.
\end{glossarylist}

\subsection*{B}
\begin{glossarylist}
\item[Backchaining] Proving an atomic goal $A$ by selecting a program clause, possibly an assumed one, whose head
  matches $A$, and then proving the clause's body. It is the only phase of search that consults the program.
  The book formalises it with focused sequents and the rules \rulename{decide}, \rulename{initial}, $\imp$L,
  $\wedge$L and $\forall$L. $\rightarrow$ \cref{ch:horn}; \book{3}, \book{44}.
\item[$\beta$-conversion, $\beta$-redex, $\beta$-reduction] A $\beta$-redex is $(\abs{x}{s})\,t$, and it reduces
  to $\subst{s}{t}{x}$. $\beta$-reduction is the computation of functional programming, and it is part of the
  equality of λProlog terms. $\rightarrow$ \cref{ch:lambda}; \book{100}--\book{101}.
\item[$\beta_0$-conversion] The special case $(\abs{x}{s})\,x = s$ of $\beta$. With $\alpha$, it is equivalent to
  $(\abs{y}{s})\,x=\subst{s}{x}{y}$ when $x$ is not free in $\abs{y}{s}$, so it only renames. It is the only $\beta$
  needed in $\Llam$. $\rightarrow$ \cref{ch:lambdatree,ch:hounif}; \book{206}, \book{220}.
\item[$\beta$-normal form] A term without $\beta$-redexes. It has the shape
  $\lambda x_1\dots\lambda x_n.\,(h\,t_1\dots t_m)$, with \emph{binder} $x_1\dots x_n$, \emph{head} $h$ and
  \emph{arguments} $t_1\dots t_m$. $\rightarrow$ \cref{ch:lambda}; \book{101}, \book{111}.
\item[$\beta\eta$-long normal form] A $\beta$-normal form whose arguments are themselves in long form and whose
  body has primitive type. Every $\beta$-normal form can be $\eta$-expanded into an equal long one; for
  $f:(i\to i)\to i$, the long form of $f$ is $\lambda w.\,f\,(\lambda u.\,w\,u)$. Higher-order unification works on
  these forms.
  $\rightarrow$ \cref{ch:lambda}; \book{111}.
\item[Big-step / small-step semantics] A big-step (natural) semantics relates a term to its final value. A
  small-step (structural operational) semantics relates a term to the result of one atomic step.
  $\rightarrow$ \cref{ch:apps}; \book{250}--\book{251}.
\item[Binary clause] A clause whose body is at most one, possibly complex, atom. Every clause can be put in this
  form by threading a continuation goal $K:\typ{o}$ (continuation-passing style).
  $\rightarrow$ \cref{ch:hoprog}; \book{130}.
\item[Binder] A syntactic construct that introduces a locally scoped name: a quantifier, $\lambda$, $\int$,
  $\Sigma$ or a local declaration. In λ-tree syntax, object-level binders are encoded by meta-level $\lambda$.
  $\rightarrow$ \cref{ch:lambdatree}.
\item[Binder mobility] The way a binder moves across levels during computation. A term-level $\lambda$ becomes a
  formula-level $\forall$ (\lp{pi}), which in turn becomes a proof-level eigenvariable. This movement lets
  λProlog programs descend under binders without managing names. $\rightarrow$ \cref{ch:lambdatree};
  \book{184}, \book{205}.
\item[Binder-merging] (Ma--Friedman) The implementation technique of renaming one side's binders to the other
  side's before comparing the bodies structurally. With name shadowing and a naive reading of levels back into
  names, it can make Church \emph{false} $\lambda a.\lambda a.a$ unify with Church \emph{true}
  $\lambda c.\lambda d.c$. $\rightarrow$ \cref{ch:lkanren}.
\item[Bisimulation, simulation] Coinductive behavioural equivalence and preorder on processes. Book Ch.~11 shows
  that a simulation written with \lp{pi} is incorrect, because \lp{pi} is intensional.
  $\rightarrow$ \cref{ch:apps}; \book{271}.
\end{glossarylist}

\subsection*{C}
\begin{glossarylist}
\item[Canonical substitution] For a problem whose equations are all flex-flex: fix one new existential
  variable $H^{\sigma}$ for each primitive type $\sigma$, and map every existential variable of type
  $\tau_1\to\dots\to\tau_n\to\sigma$ to $\lambda x_1\dots\lambda x_n.\,H^{\sigma}$. This substitution always unifies
  such a problem, which is why a flex-flex remainder counts as success. $\rightarrow$ \cref{ch:hounif};
  \book{217}--\book{218}.
\item[Capture-avoiding substitution] $\subst{s}{t}{x}$, carried out after renaming bound variables of $s$ so that
  no free variable of $t$ is captured. It is defined only up to $\alpha$. Contrast grafting.
  $\rightarrow$ \cref{ch:hh,ch:lambda}; \book{85}--\book{86}, \book{100}.
\item[Church numeral] $\ulcorner n\urcorner=\lambda f.\lambda x.\,f^n x$. The book uses Church numerals to show
  $\beta$-reduction as computation and to encode Diophantine equations as unification problems.
  $\rightarrow$ \cref{ch:lambda}; \book{102}, \book{115}.
\item[Clausal order] The nesting depth of implications to the left in a formula, analogous to the order of a type.
  $\rightarrow$ \cref{ch:hh}; \book{76}.
\item[Closed-world / open-world assumption] Whether the world, that is, the signature--program pair, can grow. Under
  the open-world assumption that λProlog adopts, generic and universal judgments coincide.
  $\rightarrow$ \cref{ch:hh}; \book{94}.
\item[Complete set of unifiers (CSU)] A set of unifiers such that every unifier is an instance of some member.
  Higher-order problems have CSUs, possibly infinite, but in general no mgu. One might ask instead for a \emph{non-redundant}
  CSU, in which no two members have a common instance, but the book notes that this cannot be guaranteed in general
  either. $\rightarrow$ \cref{ch:hounif}; \book{213}.
\item[Completeness (of a search strategy)] In the book, completeness is a property of a proof system: uniform proofs
  exist whenever proofs do. A search \emph{strategy} is a different matter. Depth-first search is incomplete.
  miniKanren's interleaving search is complete for answers, meaning every answer with a finite derivation is
  eventually produced~\cite{RozplokhasVyatkinBoulytchev2020}. It does not guarantee termination or finite
  failure. $\rightarrow$ \cref{ch:horn,ch:micro}.
\item[Computation-as-deduction] The view that computation is the construction or transformation of proofs. It
  comes in two forms: \emph{proof normalisation} (functional programming) and \emph{proof search} (logic
  programming). \emph{Computation-as-model} instead uses logic to reason about computations.
  $\rightarrow$ \cref{ch:intro}; \book{2}.
\item[\texttt{cond°}] \micro{}'s \texttt{conde}: an interleaving disjunction of conjunctions of goals. There is no
  \texttt{else} keyword. $\rightarrow$ \cref{ch:micro}.
\item[Consistent / inconsistent program] A program is inconsistent if every formula is provable from it. Every fohc
  and fohh program is consistent. Clause heads must be rigid: with the two clauses \lp{P 5 :- q.} and \lp{q.},
  whose first head is flexible, any formula becomes provable. $\rightarrow$ \cref{ch:hoprog}; \book{119}.
\item[Constant clash] The failure of first-order unification on an equation whose two sides have distinct head
  constants. $\rightarrow$ \cref{ch:intro}; \book{30}.
\item[Context (Ma--Friedman)] The set of postponed equations between static closures of unification variables. A
  solution in their algorithm is a pair (grafting, context), much like Huet's postponed flex-flex pairs.
  $\rightarrow$ \cref{ch:lkanren}.
\item[Constraint store (\micro{})] Besides the substitution \texttt{S}, a state carries three stores:
  \texttt{D} for disequalities (\texttt{≠°}), \texttt{A} for type tags (\texttt{symbol°}, \texttt{number°}) and
  \texttt{T} for absento constraints (\texttt{absent°}). After a unification, the constraints touched by its new
  bindings are checked again. $\rightarrow$ \cref{ch:micro}.
\item[Constructive negation] (Moiseenko) Negation for miniKanren that collects all the answers of $g$ and negates
  their disjunction. The states then contain universally quantified disequalities, which need a dedicated solver.
  It requires stratification, diverges on goals with infinitely many answers, and defines $\forall$ as
  $\neg\exists\neg$~\cite{Moiseenko2019}. $\rightarrow$ \cref{ch:lkanren}.
\item[Contraction] The structural rule that duplicates a hypothesis into a premise. It is the source of
  non-termination in naive provers for propositional intuitionistic logic. Contraction-free calculi such as
  Dyckhoff's terminate. $\rightarrow$ \cref{ch:apps}; \book{245}.
\item[Copy clauses] Clauses \lp{copy} that specify structural equality over an object signature, one clause per
  constructor. Adding the hypothesis \lp{copy c T} for an eigenvariable \lp{c} turns copying into object-level
  substitution, entirely inside $\Llam$. $\rightarrow$ \cref{ch:lambdatree}; \book{199}.
\item[CPS (continuation-passing style)] In book Ch.~5, the threading of a continuation goal that yields binary
  clauses. In book Ch.~10, Fischer's CPS transformation of functional programs.
  $\rightarrow$ \cref{ch:hoprog,ch:apps}; \book{131}, \book{256}.
\item[Cut rule, cut elimination] The cut rule is the lemma rule of the sequent calculus. Gentzen's theorem shows it
  is admissible~\cite{Gentzen1935}. Proof search looks only for cut-free proofs. The cut rule is unrelated to
  Prolog's \lp{!}. $\rightarrow$ \cref{ch:horn}; \book{63}, \book{72}.
\item[Cut, Prolog's \texttt{!}] The pruning operator. λProlog defines \lp{if} and \lp{not} with it, and it affects
  completeness. \micro{} has no \texttt{!}. Its closest tools are the soft cut \texttt{if°} and
  \texttt{(take° 1 g)}. $\rightarrow$ \cref{ch:hoprog,ch:micro}; \book{137}.
\end{glossarylist}

\subsection*{D}
\begin{glossarylist}
\item[D-formula, G-formula] Program clauses and goals. Each language defines them by a pair of grammars. For fohc:
  $G ::= \top\mid A\mid G\wedge G\mid G\vee G\mid \exists x\,G$ and $D ::= A\mid G\imp D\mid D\wedge D\mid\forall x\,D$.
  fohh adds $D\imp G$ and $\forall x\,G$ to the goals. $\rightarrow$ \cref{ch:horn,ch:hh}; \book{43}, \book{76}.
\item[de Bruijn index, de Bruijn level] Nameless representations of bound variables~\cite{deBruijn1972}. An
  \emph{index} counts binders outwards from the occurrence (1-based in the book). A \emph{level} counts binders
  inwards from the root. In $\lambda c.\lambda d.c$, $c$ has level 0 and index 1 when counting from 0. Teyjus uses
  indices in terms and levels in the context. ELPI uses levels throughout.
  $\rightarrow$ \cref{ch:lambdatree,ch:lkanren}; \book{192}.
\item[Decide, initial] The two rules that start and end backchaining. \rulename{decide} picks a clause $D\in\Prog$
  to focus on. \rulename{initial} closes a focused atom whose head matches the goal.
  $\rightarrow$ \cref{ch:horn}; \book{45}.
\item[\texttt{define-relation}] The \micro{} macro that defines a relation. It is inverse-$\eta$-delayed: the body
  is built only when a state arrives, and then it suspends. $\rightarrow$ \cref{ch:micro}.
\item[Definite formula] A program formula that does not contain \enquote{indefinite information of the kind let in
  by disjunctions}. Disjunction (and, by the grammar, top-level $\exists$) is banned from programs.
  $\rightarrow$ \cref{ch:horn}; \book{42}.
\item[Delay, \texttt{δ}] In \micro{}, \texttt{δ} is CHICKEN's memoising \texttt{delay}. A promise is an immature
  stream and is the unit of interleaving. $\rightarrow$ \cref{ch:micro}.
\item[Delayed (postponed) equations] Equations outside the decidable fragment, or flex-flex pairs, kept as
  constraints until further instantiation makes them tractable. Teyjus, DHKP~\cite{DowekHardinKirchnerPfenning1996}
  and the Ma--Friedman \emph{context} postpone them. ELPI raises an error by default. The book's version is
  dynamic $\Llam$. $\rightarrow$ \cref{ch:hounif,ch:lkanren}; \book{225}--\book{226}.
\item[Delayed raising] Instead of raising eagerly, keep with each existential variable the list of universals it
  may be raised over, and raise only when a binding needs it. These lists are initial segments of the prefix, so
  numerical indices represent them~\cite{NadathurLinnell2005}. $\rightarrow$ \cref{ch:hounif}; \book{225}.
\item[Determinate type] A type whose type variables are all transparent, that is, occur in the target type.
  $\rightarrow$ \cref{ch:horn}; \book{67}.
\item[Difference list, functional difference list] \lp{(dl L K)} denotes $L$ minus its suffix $K$.
  \lp{(fdl F)}, with \lp{F : list A -> list A}, denotes \lp{F nil}, and appending is function composition.
  $\rightarrow$ \cref{ch:hoprog}; \book{140}.
\item[Disagreement set] In nominal unification, the set of atoms on which two permutations differ. In Huet-style
  presentations, the multiset of pairs still to be unified. $\rightarrow$ \cref{ch:lkanren}.
\item[Disequality constraint] \texttt{≠°} in \micro{}, stored in the store \texttt{D} as the list of bindings that
  would make the two terms equal. Moiseenko's constructive negation needs universally quantified disequalities
  instead~\cite{Moiseenko2019}. $\rightarrow$ \cref{ch:micro,ch:lkanren}.
\item[Disjunction / existence property] In intuitionistic logic, $\vdash F_1\vee F_2$ implies $\vdash F_i$ for
  some $i$, and $\vdash\exists x\,F$ yields a witness. Harrop formulas guarantee these properties, which is why
  answer substitutions make sense. $\rightarrow$ \cref{ch:hh}; \book{93}.
\item[Dynamic $\Llam$] The practical regime in which $\Llam$ is not imposed statically on clauses. Instead, the
  unification problems met at run time are expected to fall inside it, as in \lp{subst R N (R N)} with \lp{R}
  and \lp{N} already instantiated. $\rightarrow$ \cref{ch:hounif}; \book{225}--\book{226}.
\end{glossarylist}

\subsection*{E}
\begin{glossarylist}
\item[E-formula] $E ::= D\mid E\wedge E\mid\exists x\,E$: a program formula with existentially quantified, local
  names. It gives modules with hidden constants a logical reading. $\rightarrow$ \book{160}--\book{162}.
\item[Eigenvariable, scoped constant] The new constant that GENERIC ($\forall$R) introduces to prove
  $\forall x\,G$. It is not in the current signature and is never instantiated. The \emph{eigenvariable
  condition} says that no logic variable created before it may be bound to a term containing it. The book's
  example is the sterile jar: to prove \lp{pi x\ bug x => in x j => dead x}, pick a new bug constant and reason
  about it alone. $\rightarrow$ \cref{ch:hh,ch:lambda}; \book{82}--\book{83}.
\item[\texttt{eigen}] (Hemann, Byrd, Friedman) The operator of miniKanren-with-symbolic-constraints that creates
  eigen objects and installs an \enquote{eigen-absento} constraint over the list of all variables created so far.
  It is supported only together with \texttt{==}~\cite{HemannByrdFriedman2013}. $\rightarrow$ \cref{ch:lkanren}.
\item[Elementary substitution] Huet's partial bindings for a flex-rigid pair: imitation and projection. In
  Ma--Friedman's reading, each unification variable takes every visible bound variable as an argument.
  $\rightarrow$ \cref{ch:hounif}.
\item[ELPI] An embeddable λProlog interpreter~\cite{DunchevEtAl2015}. It uses de Bruijn levels and annotates
  unification variables with levels $X^j$. It reports an error outside the pattern fragment by default, and it
  offers syntactic constraints. $\rightarrow$ \cref{ch:lkanren}.
\item[Empty type] A type with no closed inhabitant. Higher-order intuitionistic logic does not exclude empty types,
  and they matter for $\exists$ answers and for unification \emph{solutions}. $\rightarrow$ \cref{ch:lambda};
  \book{108}.
\item[\texttt{eq?}-invariant] (\micro{}) A unification that adds no binding returns the very same state object.
  The disequality and entailment checks rely on this to recognise that nothing new was learned.
  $\rightarrow$ \cref{ch:micro}.
\item[Essentially universal / essentially existential] A classification of bound-variable occurrences. In a
  goal, an occurrence is essentially universal if it is bound by a positive $\forall$, a negative $\exists$, or a
  term-level $\lambda$. Otherwise it is essentially existential. Clauses use the dual classification. During search,
  essentially existential variables become logic variables and essentially universal ones become eigenvariables.
  $\rightarrow$ \cref{ch:lambdatree}; \book{205}.
\item[$\eta$-conversion] $\abs{x}{(t\,x)} = t$ when $x$ is not free in $t$. $\rightarrow$ \cref{ch:lambda};
  \book{101}.
\item[Evaluation context] A term with a hole that marks where the next redex is. In book Ch.~10 it is a meta-level
  function \lp{tm -> tm}. $\rightarrow$ \cref{ch:apps}; \book{253}.
\item[Ex falso quodlibet] From $\bot$, infer anything. Intuitionistic logic is minimal logic plus this rule.
  $\rightarrow$ \cref{ch:hh}; \book{90}.
\item[Excluded middle] $B\vee\neg B$: valid classically, rejected intuitionistically. It is one of the
  counterexamples that motivate banning disjunction from programs. $\rightarrow$ \cref{ch:horn}; \book{42}.
\item[Existential generalisation] From $M(t)$ infer $\exists x\,M(x)$. It explains the multiple answers of
  \lp{subst} when run backwards. $\rightarrow$ \cref{ch:lambdatree}; \book{203}.
\item[Explicit substitution] Calculi that make substitution a term constructor, so it can be delayed and composed.
  Examples are $\lambda\sigma$ (DHKP) and Teyjus's suspension notation. $\rightarrow$ \cref{ch:lkanren};
  \book{209}.
\end{glossarylist}

\subsection*{F}
\begin{glossarylist}
\item[Flexible term, rigid term] A $\beta\eta$-normal term is \emph{flexible} if its head is an existential
  (unification) variable. It is \emph{rigid} if its head is a constant or a universally bound or $\lambda$-bound
  variable. The names are suggestive: Ma and Friedman, following Huet, explain that a flexible term is one that
  \enquote{can be removed through a substitution (or a grafting)}, while a substitution never changes the head of
  a rigid term~\cite{MaFriedman2021}. $\rightarrow$ \cref{ch:lambda,ch:hounif}; \book{111}, \book{214}.
\item[Flexible atom, rigid atom, flexible goal] An atom is flexible if its head is a variable and rigid if its head
  is a non-logical constant. Clause heads must be rigid. A flexible goal whose head is still uninstantiated when it
  is called is underconstrained. $\rightarrow$ \cref{ch:hoprog}; \book{119}, \book{132}.
\item[Flex-flex, flex-rigid, rigid-rigid] The classification of equations by the heads of their two sides.
  Rigid-rigid pairs simplify by decomposition or fail by a clash. Flex-rigid pairs call for imitation or
  projection. Flex-flex pairs always have unifiers, though not necessarily solutions, and are postponed.
  $\rightarrow$ \cref{ch:hounif}; \book{113}, \book{217}.
\item[Focused proof] Andreoli's generalisation of the two-phase structure of goal reduction followed by
  backchaining. Focusing groups small rules into large ones~\cite{Andreoli1992}. $\rightarrow$ \cref{ch:horn};
  \book{73}.
\item[fohc, fohh, hohc, hohh, hohh$^{+}$] The book's languages: first-order Horn clauses \book{43}, first-order
  hereditary Harrop formulas \book{76}, higher-order Horn clauses \book{122}, higher-order hereditary Harrop
  formulas \book{124}, and hohh$^+$, which additionally allows atomic clauses headed by an essentially universal
  variable \book{125}. $\rightarrow$ \cref{ch:horn,ch:hh,ch:hoprog}.
\item[Free for] $t$ is free for $x$ in $B$ if substituting $t$ for $x$ in $B$ captures no free variable of $t$.
  $\rightarrow$ \cref{ch:hh}; \book{85}.
\item[Fresh, \texttt{fresh°}] miniKanren's existential quantifier, which introduces logic variables. In \micro{},
  \texttt{fresh°} introduces them one at a time, and each allocation suspends once. For the αKanren meaning, see
  \emph{Atom (nominal)}. $\rightarrow$ \cref{ch:micro}.
\item[Freshness constraint] $a\mathbin{\#}t$: the atom $a$ does not occur free in $t$, also after later
  instantiation of the variables in $t$. It is written \texttt{hash} in αKanren~\cite{UrbanPittsGabbay2004}.
  $\rightarrow$ \cref{ch:lkanren}.
\end{glossarylist}

\subsection*{G}
\begin{glossarylist}
\item[Generic judgment, universal judgment] A generic judgment is parametric in a new object and is proved once,
  with an eigenvariable. A universal judgment holds for all terms. The two coincide under the open-world
  assumption. They correspond to the intensional and the extensional readings of $\forall$.
  $\rightarrow$ \cref{ch:hh}; \book{83}, \book{94}.
\item[GENERIC] The rule for $\forall x\,G$: add a new constant $c$ to the signature and prove $\subst{G}{c}{x}$.
  It is \lp{pi x\ G} in λProlog and \texttt{all} in λKanren.
  \[\inferrule{\seq{\Sig\cup\{c\}}{\Prog}{\subst{G}{c}{x}}}{\seq{\Sig}{\Prog}{\forall x\,G}}\;\rulename{generic}\quad(c\notin\Sig)\]
  $\rightarrow$ \cref{ch:hh}; \book{39}, \book{82}.
\item[Goal (miniKanren)] In \micro{}, a procedure from a state to a stream of states. \texttt{✓°} succeeds once
  and \texttt{✗°} fails. $\rightarrow$ \cref{ch:micro}.
\item[Goal-directed search, goal reduction] A non-atomic goal is decomposed by its top-level connective,
  independently of the program. This fixes the \enquote{search semantics} of the connectives. Only atomic goals
  consult the program. $\rightarrow$ \cref{ch:intro,ch:horn}; \book{3}, \book{39}.
\item[Grafting] First-order, capture-permitting replacement of a metavariable by a term, written $\{X\mapsto t\}$:
  grafting $X\mapsto a$ into $\lambda a.X$ gives $\lambda a.a$. DHKP \enquote{pre-cook} terms so that grafting
  respects scopes. Ma and Friedman combine grafting with de Bruijn levels~\cite{MaFriedman2021}.
  $\rightarrow$ \cref{ch:lkanren}.
\item[Grounding] In \micro{}, \texttt{°->list/ground} evaluates a reified answer, applying its $\lambda$ to its own
  variable names. Constraints are dropped. $\rightarrow$ \cref{ch:micro}.
\end{glossarylist}

\subsection*{H}
\begin{glossarylist}
\item[Harrop formula] $H ::= \top\mid A\mid H\wedge H\mid B\imp H\mid\forall x\,H$. These formulas guarantee the
  disjunction and existence properties at the root of a proof. $\rightarrow$ \cref{ch:hh}; \book{93}.
\item[Herbrand universe $H_1^{\Sig}$, $H_2^{\Sig}$] The terms allowed as quantifier instances. $H_1^{\Sig}$ (for
  hohc) contains the $\lambda$-normal terms over $\Sig$ without $\forall$ and $\imp$. $H_2^{\Sig}$ (for hohh)
  contains those without $\imp$. $\rightarrow$ \cref{ch:hoprog}; \book{122}.
\item[Hereditary Harrop formula] A D-formula of fohh or hohh. Equivalently, a formula with no positive occurrence
  of a disjunction or an existential. Unlike Horn clauses, these formulas let the signature and the program grow
  during search. $\rightarrow$ \cref{ch:hh}; \book{4}, \book{76}.
\item[Higher-order abstract syntax (HOAS)] The principle of encoding object-level binders by meta-level
  binders~\cite{PfenningElliott1988}. It is ambiguous between a functional reading, in which binders are host
  functions that cannot be analysed, and the logical reading, λ-tree syntax. $\rightarrow$ \cref{ch:lambdatree};
  \book{204}.
\item[Higher-order logic] The book distinguishes three senses: second-order model theory; proof theory with
  predicate quantification and comprehension; and Church's simple theory of types, with $\lambda$-terms and
  quantification at higher types. The book works in the third. $\rightarrow$ \cref{ch:intro}; \book{5}.
\item[Higher-order pattern] A term in which every occurrence of an existential variable $F$ is applied to distinct
  variables, either essentially universal or $\lambda$-bound, whose binders lie within the scope of $F$. For
  example, $\abs{x}{\abs{y}{F\,y\,x}}$ is a pattern, while $\abs{x}{F\,x\,x}$ (repeated argument) and $F\,c$ with $c$
  a signature constant (not a bound variable) are not.
  $\rightarrow$ \cref{ch:lambdatree,ch:hounif}; \book{205}, \book{220}.
\item[Higher-order pattern unification ($\Llam$-unification)] Unification restricted to patterns. It is decidable,
  and unifiable problems have an mgu~\cite{Miller1991,Nipkow1993}. Its steps include an occurs check and variable
  elimination with pruning. $\rightarrow$ \cref{ch:hounif}; \book{220}--\book{223}.
\item[Higher-order programming] Passing predicates as arguments, as in \lp{mappred} and \lp{foldl}.
  $\rightarrow$ \cref{ch:hoprog}; \book{126}.
\item[Higher-order unification] Unification of simply typed $\lambda$-terms modulo $\lambda$-conversion. It is
  undecidable (the book reduces Post's correspondence problem and Hilbert's tenth problem to it), has no mgus in
  general, and is semi-decided by Huet's procedure~\cite{Huet1975}. $\rightarrow$ \cref{ch:lambda,ch:hounif};
  \book{111}, \book{114}--\book{115}, \book{214}.
\item[Horn clauses] The class of programs in which the signature and the program stay fixed during search: they are
  \enquote{flat and global}. $\rightarrow$ \cref{ch:horn}; \book{43}, \book{46}.
\item[Huet's procedure, pre-unification] Repeated simplification, imitation and projection, stopping when only
  flex-flex pairs remain~\cite{Huet1975}. The result is a pre-unifier, whose search space is explored as a
  matching tree. $\rightarrow$ \cref{ch:hounif}; \book{218}--\book{219}.
\item[Hypothetical reasoning] Proving a goal under temporarily assumed clauses, by AUGMENT. The book's example is
  the question \enquote{If person X took just one more course Y, will he or she graduate with a CS degree?}, posed as the query \lp{fact (finished X Y) => (fact (graduates X), fact (cs_major X))}.
  $\rightarrow$ \cref{ch:hh}; \book{79}--\book{80}.
\end{glossarylist}

\subsection*{I}
\begin{glossarylist}
\item[\texttt{if}, \texttt{not}] λProlog's conditional and negation-as-failure, both defined with the cut.
  $\rightarrow$ \cref{ch:hoprog}; \book{137}.
\item[\texttt{if°} (soft cut)] \texttt{(if° g? gt gf)}: if \texttt{g?} has at least one answer, every one of
  its answers continues with \texttt{gt}; otherwise \texttt{gf} runs. It corresponds to miniKanren's
  \texttt{conda} and LogicT's \texttt{ifte}. λProlog's \lp{if} becomes \texttt{(if° (take° 1 P) Q R)}.
  $\rightarrow$ \cref{ch:micro}.
\item[Imitation, projection] Huet's partial bindings for a flex-rigid pair. \emph{Imitation} copies the rigid head:
  $F:=\lambda\bar x.\,c\,(H_1\bar x)\cdots(H_m\bar x)$. \emph{Projection} selects an argument:
  $F:=\lambda\bar x.\,x_i\,(H_1\bar x)\cdots(H_k\bar x)$, which requires matching target types.
  $\rightarrow$ \cref{ch:hounif}; \book{215}.
\item[Intensional / extensional equality] Intensional equality compares the structure of $\lambda$-terms.
  Extensional equality compares what they denote. λProlog's \lp{=} is intensional (modulo $\lambda$-conversion).
  $\rightarrow$ \cref{ch:hoprog}; \book{146}.
\item[Intensional / extensional $\forall$] \lp{pi} is intensional: it requires one generic proof, so it cannot
  express \enquote{for all instances}. This is why book Ch.~11's simulation is incorrect. Jin et al.\ implement an
  extensional $\forall$ in miniKanren, which is a different logic~\cite{JinEtAl2021}.
  $\rightarrow$ \cref{ch:hh,ch:apps}; \book{83}, \book{94}.
\item[Interleaving] miniKanren's search discipline. In \micro{}, \texttt{orª} uses
  \texttt{append§/interleaved/2}, which swaps the two streams whenever the first one is a suspension, and
  conjunction (\texttt{append-map§}) interleaves too. The right-nested \texttt{or°} gives the first of three
  infinite branches about half of the answers. $\rightarrow$ \cref{ch:micro}.
\item[Inverse-$\eta$ delay] The wrapping \texttt{(λ (s) (δ ((and° g …) s)))} that \texttt{define-relation} puts
  around a body, so that a recursive relation without \texttt{fresh°} does not loop while the goal tree is built.
  $\rightarrow$ \cref{ch:micro}.
\item[Inversion of a pattern substitution] The core step of pattern unification. To solve $X\,\bar y = t$, where
  $\bar y$ are distinct bound variables, abstract $t$ over $\bar y$, that is, $X:=\lambda\bar y.\,t$. This works only
  if every bound variable free in $t$ is among $\bar y$ and $X$ does not occur in $t$; otherwise pruning or the occurs
  check takes over. DHKP write the inverse as $b[\xi]^{-1}$, and Nipkow's code remaps de Bruijn indices.
  $\rightarrow$ \cref{ch:hounif}.
\item[Invertible rule] A rule whose premises are provable if and only if its conclusion is. Such rules can be
  applied eagerly without losing completeness. $\rightarrow$ \cref{ch:apps}; \book{243}.
\item[Iterative deepening] Retrying a search with bound $N+1$ after failing with bound $N$. In book Ch.~9 it bounds
  the multiplicity of $\exists$-formulas. $\rightarrow$ \cref{ch:apps}; \book{237}, \book{239}.
\end{glossarylist}

\subsection*{K}
\begin{glossarylist}
\item[Kind] The \enquote{type} of a type constructor, of the form \lp{type -> ... -> type}, introduced by a
  \lp{kind} declaration. $\rightarrow$ \cref{ch:intro}; \book{11}--\book{12}.
\item[Kripke model] The possible-worlds semantics of intuitionistic logic. Its worlds correspond to
  signature--program pairs. $\rightarrow$ \cref{ch:hh}; \book{94}.
\end{glossarylist}

\subsection*{L}
\begin{glossarylist}
\item[$\Llam$] The sublanguage of hohh in which quantification is never over predicate types, and every
  essentially existential variable is applied only to distinct essentially universal variables bound within its
  scope. Unification inside $\Llam$ needs only $\beta_0$, is decidable and produces mgus.
  $\rightarrow$ \cref{ch:lambdatree,ch:hounif}; \book{205}--\book{206}, \book{220}.
\item[$\lambda$-conversion, $\lambda$-normal form] $\lambda$-conversion is the closure of $\alpha$, $\beta$ and
  $\eta$, and it is the equality of the logic. The $\lambda$-normal form $\lambda norm(s)$ is the $\beta\eta$-normal
  form. $\rightarrow$ \cref{ch:lambda}; \book{100}--\book{102}.
\item[λ-tree syntax] HOAS with simply typed $\lambda$-terms, taken modulo $\alpha$, a weak $\beta$ ($\beta_0$) and
  possibly $\eta$, used as a data structure that unification can still analyse. $\rightarrow$
  \cref{ch:lambdatree}; \book{204}.
\item[λKanren] miniKanren extended with higher-order logic programming. The 2020 version has terms \texttt{tie}
  and \texttt{app}, goals \texttt{assume} ($D\imp G$) and \texttt{all} ($\forall x\,G$), and Miller's pattern
  unification over an explicit prefix~\cite{MaLuFriedman2020}. The 2021 paper proposes static
  closures~\cite{MaFriedman2021}. $\rightarrow$ \cref{ch:lkanren}.
\item[Left- and right-introduction rules] Right rules introduce a connective on the right of $\longrightarrow$ and
  implement goal reduction. Left rules decompose the clause in focus and implement backchaining.
  $\rightarrow$ \cref{ch:horn}; \book{40}, \book{42}, \book{45}.
\item[Legitimate substitution] For a problem $\forall\bar u\,\exists\bar x\,\forall\bar w\,[\dots]$, a substitution
  for $\bar x$ in whose range no $w_j$ occurs free. Only legitimate substitutions can be unifiers. This is the
  eigenvariable condition, stated for unification. $\rightarrow$ \cref{ch:lambda}; \book{109}.
\item[Leibniz equality] $x=y$ if and only if $\forall P\,(P\,x\imp P\,y)$. $\rightarrow$ \cref{ch:hoprog};
  \book{120}.
\item[Level of a unification variable] The scope depth of a variable. ELPI's $X^j$ may mention only context names
  below $j$. In Teyjus, a variable may be bound only to terms whose constants have a tag no larger than its own.
  Binding a variable must lower the levels of the variables in the bound term.
  $\rightarrow$ \cref{ch:hounif,ch:lkanren}.
\item[Locally nameless] A representation with bound variables as de Bruijn indices and free variables as
  names~\cite{McBrideMcKinna2004,Chargueraud2012}. Eigen constants would play the role of the names.
  $\rightarrow$ \cref{ch:lkanren}.
\item[Logic variable] A placeholder that search instantiates by unification. It is the operational device behind
  the delayed choice of terms in the $\exists$R, $\forall$L and \rulename{initial} rules.
  $\rightarrow$ \cref{ch:horn}; \book{61}.
\end{glossarylist}

\subsection*{M}
\begin{glossarylist}
\item[Matching] A unification problem in which only one side contains existential variables. Higher-order matching
  is decidable. $\rightarrow$ \cref{ch:hounif}; \book{227}.
\item[Matching tree] The search tree of Huet's procedure. Its nodes are multisets of equations and its arcs are
  substitutions. A failure node is a problem that obviously has no solution (for example, a rigid-rigid clash); a success node
  contains at most flex-flex pairs, and reaching one is a pre-unification. Some paths may be infinite.
  $\rightarrow$ \cref{ch:hounif}; \book{219}.
\item[Matrix] The conjunction of equations $t_1=s_1\wedge\dots\wedge t_m=s_m$ under the quantifier prefix of a
  unification problem. $\rightarrow$ \cref{ch:hounif}; \book{211}.
\item[May judgment, must judgment] A may judgment asserts that some computation path exists. A must judgment states
  a property of all paths, and needs negation-as-failure or a meta-level $\forall$.
  $\rightarrow$ \cref{ch:apps}; \book{269}.
\item[Metalanguage, object language] λProlog, the language in which we write, and the language encoded in it.
  $\rightarrow$ \cref{ch:intro}; \book{22}.
\item[Minimal logic] Logic in which $\bot$ is an ordinary constant. Negation is defined as
  $\neg B := B\imp\cst{false}$. $\rightarrow$ \cref{ch:hh}; \book{89}--\book{90}.
\item[miniKanren] A family of small relational languages embedded in a host language~\cite{FriedmanEtAl2018,Byrd2009}.
  $\rightarrow$ \cref{ch:micro}.
\item[Mixed prefix] The quantifier prefix $\forall\exists\forall\dots$ of a unification problem. In a
  $\forall\exists\forall$ problem the outer universals act as the signature and the inner ones as eigenvariables,
  and no raising is needed. $\rightarrow$ \cref{ch:lambda,ch:hounif}; \book{105}, \book{211}.
\item[Module, signature (Teyjus)] A module is a named collection of declarations and clauses. Its signature
  restricts what is visible from outside, and hidden constants are read as existentially quantified (book
  Ch.~6). $\rightarrow$ \cref{ch:intro}; \book{152}--\book{153}, \book{277}.
\item[Most general unifier (mgu)] A unifier $\theta$ such that every other unifier is $\theta$ followed by a
  further substitution. In the first-order case an mgu exists whenever a unifier does, and it is unique up to
  renaming. Pattern problems have mgus too, but general higher-order problems do not.
  $\rightarrow$ \cref{ch:intro,ch:hounif}; \book{28}, \book{213}, \book{222}.
\item[μKanren] Hemann and Friedman's minimal miniKanren, in which goals are functions from a state to a stream of
  states~\cite{HemannFriedman2013}. \micro{} is a μKanren with constraint stores. $\rightarrow$ \cref{ch:micro}.
\item[Multiplicity] In the classical prover of book Ch.~9, the number of times an $\exists$-formula may be
  instantiated. Bounding it, and raising the bound after each failure, gives a complete semi-decision procedure.
  $\rightarrow$ \cref{ch:apps}; \book{237}.
\end{glossarylist}

\subsection*{N}
\begin{glossarylist}
\item[$\nabla$ (nabla)] The Miller--Tiu generic quantifier for fresh names, in a logic with definitions (closed
  world). It is not part of λProlog. $\rightarrow$ \cref{ch:hh,ch:apps}; \book{94}, \book{275}.
\item[Name shadowing] An inner binder that reuses an outer binder's name, so that the outer binder can no longer be
  named. $\rightarrow$ \cref{ch:lkanren}.
\item[Natural deduction] Introduction and elimination rules, with hypotheses discharged by $\imp$I, $\vee$E and
  $\exists$E. Book Ch.~9 encodes it with \lp{=>} and \lp{pi}. $\rightarrow$ \cref{ch:apps}; \book{231}.
\item[Negation-as-failure] \lp{not G} succeeds if and only if \lp{G} finitely fails. It is non-logical and unsound
  on non-ground goals. In \micro{} it is written \texttt{(if° (take° 1 g) ✗° ✓°)}. Without the \texttt{take° 1},
  it diverges whenever the search for further answers of \texttt{g} does not terminate, even if \texttt{g} has a
  single answer (\cref{sec:glossary-bridge} runs an example). $\rightarrow$ \cref{ch:hoprog,ch:apps}; \book{137},
  \book{269}.
\item[Negation normal form, prenex normal form] NNF has no implications and applies negation only to atoms. A
  prenex form has all its quantifiers in front. $\rightarrow$ \cref{ch:lambdatree}; \book{195}.
\item[Nominal logic, nominal unification] A first-order treatment of binding with names, swappings and freshness.
  Nominal unification works modulo $\alpha$ and never generates new names~\cite{UrbanPittsGabbay2004}.
  $\rightarrow$ \cref{ch:lkanren}; \book{210}.
\end{glossarylist}

\subsection*{O}
\begin{glossarylist}
\item[\texttt{o}] The type of propositions, named after Church's omicron. The quantifier names \lp{pi} and
  \lp{sigma} likewise follow Church's $\Pi$ and $\Sigma$. $\rightarrow$ \cref{ch:intro}; \book{74}.
\item[Object-level substitution] Substituting into the body of an encoded binder. In λ-tree syntax it is meta-level
  $\beta$-reduction, as in the clause \lp{subst R N (R N).} with \lp{R : tm -> tm}. That clause lies outside
  $\Llam$, but it is used only with \lp{R} and \lp{N} instantiated (dynamic $\Llam$). Copy clauses specify the same
  relation inside $\Llam$. $\rightarrow$ \cref{ch:lambdatree}; \book{180}, \book{199}, \book{225}.
\item[O-proof] A derivation that uses only goal reduction (right rules) followed by backchaining (focused left
  rules). Every provable fohc sequent has an O-proof. $\rightarrow$ \cref{ch:horn}; \book{45}.
\item[Occurs check] $x=t$ fails when $x$ occurs properly in $t$. In pattern unification, the corresponding failure
  is a flexible head occurring rigidly on the other side. \micro{} performs the check when it extends the
  substitution. $\rightarrow$ \cref{ch:intro,ch:hounif}; \book{30}, \book{223}.
\item[Order of a type] $0$ for non-functional types, and $\max(\mathrm{ord}\,\tau_1+1,\mathrm{ord}\,\tau_2)$ for
  $\tau_1\to\tau_2$. $\rightarrow$ \cref{ch:intro,ch:lambda}; \book{13}.
\end{glossarylist}

\subsection*{P}
\begin{glossarylist}
\item[Pattern] See \emph{Higher-order pattern}. Ma and Friedman's phrasing, for an application form $F\,a$:
  \enquote{if $F$ is a unification variable, then $a$ must be a list of zero or more distinct bound variables}.
\item[Peirce's formula] $((p\imp q)\imp p)\imp p$: valid classically but not intuitionistically. It has no
  O-proof. $\rightarrow$ \cref{ch:hh}; \book{89}.
\item[Pervasive] Built in and globally available: sorts such as \lp{int}, constants such as \lp{nil} and \lp{::},
  and predicates such as \lp{is}. $\rightarrow$ \cref{ch:intro}; \book{14}, \book{49}.
\item[\texttt{pi}, \texttt{sigma}] λProlog's quantifiers, constants of type \lp{(A -> o) -> o}. $\forall x\,G$
  is written \lp{pi x\ G}, the application of \lp{pi} to the abstraction \lp{x\ G}, so the term-level $\lambda$
  is the only binder. $\rightarrow$ \cref{ch:intro,ch:lambda}; \book{36}, \book{99}.
\item[$\pi$-calculus] Milner's calculus of mobile processes, in which names are the only data. Book Ch.~11 encodes
  it in λ-tree syntax, using \lp{in} and \lp{nu} as binders. $\rightarrow$ \cref{ch:apps}; \book{261}.
\item[Polarity, positive and negative occurrence] An occurrence is positive if it lies to the left of an even
  number of implications and negative if odd. Cut-free search preserves polarity, unless implications appear
  inside atoms. $\rightarrow$ \cref{ch:hh,ch:hoprog}; \book{76}, \book{121}--\book{122}.
\item[Pre-cooking] (DHKP) Translating a metavariable that occurs under $n$ binders into $X[\uparrow^n]$, so that
  first-order grafting agrees with capture-avoiding substitution~\cite{DowekHardinKirchnerPfenning1996}.
  $\rightarrow$ \cref{ch:lkanren}.
\item[Pre-unifier] A substitution that reduces a problem to flex-flex pairs only. Such a remainder always has a
  unifier, the canonical substitution. $\rightarrow$ \cref{ch:hounif}; \book{217}--\book{218}.
\item[\texttt{project°}] \micro{}'s escape to Scheme: it binds Scheme variables to deep-walked values, in which
  unbound variables remain variable records, and then runs goals. It is non-relational.
  $\rightarrow$ \cref{ch:micro}.
\item[Proof term] A $\lambda$-term of type \lp{proof} that records a derivation, as in the judgment \lp{P # A}
  (\enquote{$P$ is a proof of $A$}). $\rightarrow$ \cref{ch:apps}; \book{233}.
\item[Pruning] Removing from a variable its dependency on arguments that cannot occur in any solution:
  $G:=\lambda\bar d.\,(G'\,\bar e)$, where $\bar e$ keeps only the allowed arguments. $\rightarrow$
  \cref{ch:hounif}; \book{223}.
\item[Purification] The step of \micro{}'s reification that drops constraints mentioning variables unreachable from
  the answer. $\rightarrow$ \cref{ch:micro}.
\end{glossarylist}

\subsection*{R}
\begin{glossarylist}
\item[Raising] $\forall y\,\exists x\,G \rightsquigarrow \exists h\,\forall y\,\subst{G}{(h\,y)}{x}$: an existential in the
  scope of a universal becomes the application of a new, higher-type existential quantified further out. Raising
  preserves $\Llam$. The book calls it a dual of Skolemization: Skolemization introduces a new \emph{constant} of a
  raised type, raising a new \emph{existential variable} of a raised type. $\rightarrow$ \cref{ch:lambda,ch:hounif};
  \book{108}, \book{208}, \book{212}.
\item[Reduction rules] AND, OR, INSTAN, AUGMENT, GENERIC and TRUE: the search meaning of each logical symbol.
  $\rightarrow$ \cref{ch:horn}; \book{39}.
\item[Reduction-free fragment $L^{\beta}_{\lambda}$] (ELPI) The patterns in which a unification variable $X^j$ is
  applied exactly to the bound variables $x_j\dots x_{j+k-1}$ that follow its level. Binding such a variable needs no
  $\beta$-reduction and takes constant time, but the fragment is not stable under resolution~\cite{DunchevEtAl2015}.
  $\rightarrow$ \cref{ch:lkanren}.
\item[Reification] Turning the final state into a presentable answer. \micro{} produces a closed Scheme expression
  \texttt{(λ (α β …) constraint-form … term)}, naming free variables in order of first occurrence. αKanren prints
  noms as \texttt{a.0}, \texttt{a.1}, \dots{} and adds freshness constraints. $\rightarrow$ \cref{ch:micro}.
\item[Relation (\micro{})] A procedure from terms to a goal, conventionally named \texttt{name°} and defined with
  \texttt{define-relation}. $\rightarrow$ \cref{ch:micro}.
\item[Relative complement] (Jin et al.) A goal that describes the values of a variable not covered by a state. It
  is used to implement an extensional $\forall$~\cite{JinEtAl2021}. $\rightarrow$ \cref{ch:lkanren}.
\item[Representation independence] Replacing a module by an equivalent one does not change the meaning of the
  larger program. $\rightarrow$ \book{151}.
\item[Rigid] See \emph{Flexible term, rigid term} and \emph{Flexible atom}.
\end{glossarylist}

\subsection*{S}
\begin{glossarylist}
\item[Scope check, eigenvariable condition] No variable older than an eigenvariable may reach it. The
  implementations differ: an explicit quantifier prefix (Miller; λKanren 2020), universe indices or levels
  (Teyjus, ELPI), an absento-like constraint over the variables in scope (\texttt{eigen}), or an end-of-scope check,
  which is fragile. $\rightarrow$ \cref{ch:lkanren,ch:micro}.
\item[Scope extrusion] The book's name for the classical equivalence
  $B_1\vee(B_2\imp B_3)\equiv B_2\imp(B_1\vee B_3)$, which moves the scope of an implication over a disjunction. It
  fails intuitionistically: $p\vee(p\imp q)$ has no O-proof, although its classical equivalent $p\imp(p\vee q)$ does. The
  term is borrowed from the $\pi$-calculus, where it means sending a restricted name out and re-closing its scope at
  the receiver. $\rightarrow$ \cref{ch:hh,ch:apps}; \book{89}, \book{95}.
\item[Scope restriction] Ma and Friedman's name for the eigenvariable condition in unification: a $\forall$-bound
  variable may occur in the instantiation of an $\exists$-variable only if it precedes that variable in the
  prefix. $\rightarrow$ \cref{ch:lkanren}.
\item[Shallow embedding] Implementing a logic language as host-language procedures and macros, as miniKanren
  does in Scheme, instead of writing an interpreter or compiler for it. λKanren keeps this
  style~\cite{MaLuFriedman2020}. $\rightarrow$ \cref{ch:lkanren,ch:micro}.
\item[Sequent] $\seq{\Sig}{\Prog}{G}$: signature, program and goal. It is read both as the state of an interpreter
  and as a logical judgment. $\rightarrow$ \cref{ch:horn}; \book{39}.
\item[Signature] A set of typed non-logical constants. In Teyjus, a signature is also a module's interface
  (\lp{sig}). $\rightarrow$ \cref{ch:intro}; \book{17}, \book{153}.
\item[Signature--program pair] $\langle\Sig,\Prog\rangle$, the context in which queries are posed. Under the
  open-world assumption these pairs are the worlds of a Kripke model. $\rightarrow$ \cref{ch:horn,ch:hh};
  \book{39}, \book{94}.
\item[Simplification] Repeated decomposition, or clash, of rigid-rigid pairs in higher-order unification.
  $\rightarrow$ \cref{ch:hounif}; \book{214}.
\item[Skolemization] Replacing an existential in an assumption by a function of the enclosing universals, as in
  $\forall x\,\subst{D}{f\,x}{y}$. It is the dual of raising. $\rightarrow$ \cref{ch:lambda}; \book{108}.
\item[Soft cut] See \texttt{if°}. A soft cut is not complete in general.
\item[Solution] A unifier whose range has free variables only among the outer universals. A solution amounts to a
  proof of the quantified equation. $\rightarrow$ \cref{ch:lambda,ch:hounif}; \book{109}.
\item[Solved form] A set of equations $x=t$ in which each left-hand variable occurs nowhere else. It is read as an
  idempotent mgu. $\rightarrow$ \cref{ch:intro}; \book{29}.
\item[Static closure] $\langle\phi;t\rangle$ (Ma--Friedman): a term paired with the scope, the list of binder names
  around it, manipulated by $\mathrm{ext}$ and $\mathrm{db}$. The authors call the paper \enquote{a preview} and
  leave soundness and completeness to future work~\cite{MaFriedman2021}. $\rightarrow$ \cref{ch:lkanren}.
\item[Stream (\texttt{§})] In \aux{}, a stream is \texttt{'()}, a mature pair \texttt{(state . §)}, or an immature
  promise. $\rightarrow$ \cref{ch:micro}.
\item[Subject reduction] Evaluation preserves types. $\rightarrow$ \cref{ch:lambdatree}; \book{191}.
\item[Substitution (\micro{})] The store \texttt{S}: a triangular substitution kept in a skew-binary random-access
  list indexed by variable number. $\rightarrow$ \cref{ch:micro}.
\item[Subsumption (\micro{})] Dropping a constraint that another constraint implies, or that a type tag or an
  absento constraint makes vacuous. $\rightarrow$ \cref{ch:micro}.
\item[Suspension] Four different things, depending on the source; see the pitfall at the start of this appendix.
\item[Swapping] $(a\;b)$ exchanges the atoms $a$ and $b$ throughout a term. Permutations are lists of swappings.
  $\rightarrow$ \cref{ch:lkanren}.
\item[Syntactic constraint] (ELPI) A goal suspended on some variables with \texttt{declare\_constraint} and resumed
  when one of them is assigned. It keeps its hypothetical context. $\rightarrow$ \cref{ch:lkanren}.
\end{glossarylist}

\subsection*{T}
\begin{glossarylist}
\item[\texttt{take°}] \texttt{(take° n g)} keeps the first \texttt{n} answers of \texttt{g}; a negative
  \texttt{n} keeps them all. It is a goal, so it nests, and \texttt{(take° 1 g)} is miniKanren's \texttt{once}.
  \texttt{μkanren-run} is \texttt{take°} around a \texttt{fresh°} query. $\rightarrow$ \cref{ch:micro}.
\item[Tactic, tactical] A tactic is a relation \lp{goal -> goal -> o} from a primitive goal to the goal of its
  premises. Tacticals such as \lp{then}, \lp{orelse}, \lp{repeat}, \lp{try} and \lp{maptac} combine tactics. Book
  Ch.~9 recasts the LCF architecture relationally. $\rightarrow$ \cref{ch:apps}; \book{241}, \book{243}.
\item[Target type] $\tau_0$ in $\tau_1\to\dots\to\tau_n\to\tau_0$. $\rightarrow$ \cref{ch:intro}; \book{13}.
\item[Term reduction, variable elimination] Two of the three transformations of first-order unification; the third
  is reorientation. Term reduction decomposes equal heads. Variable elimination replaces a variable by its value
  everywhere. $\rightarrow$ \cref{ch:intro}; \book{28}.
\item[Teyjus] The λProlog implementation used by the book, with compiler \texttt{tjcc}, linker \texttt{tjlink} and
  simulator \texttt{tjsim}. Teyjus~2 uses universe-level tags. $\rightarrow$ \cref{ch:intro}; \book{277}.
\item[Trace, complete trace, trace equivalence] A trace is a sequence of actions of a process, and a complete trace
  ends in a terminal process. Two processes are trace equivalent if no trace separates them.
  $\rightarrow$ \cref{ch:apps}; \book{267}, \book{269}, \book{271}.
\item[Transparent type variable] A type variable that occurs in the target type. $\rightarrow$ \cref{ch:horn};
  \book{67}.
\item[Type lowering] A polarity-indexed translation from a relation at a higher type to formulas over base-type
  relations. It generates copy clauses. $\rightarrow$ \cref{ch:lambdatree}; \book{201}--\book{202}.
\end{glossarylist}

\subsection*{U}
\begin{glossarylist}
\item[Unification problem, unifier] In the first-order case, a problem is a finite multiset of equations, and a
  unifier is a type-preserving substitution that makes both sides of each equation equal. In general, a problem
  is a quantified formula $Q_1x_1\dots Q_nx_n[\bigwedge_j t_j=s_j]$, and a unifier is a legitimate substitution
  (no inner universals in its range) that makes each pair $\lambda$-convertible.
  $\rightarrow$ \cref{ch:intro,ch:lambda}; \book{26}, \book{105}, \book{109}.
\item[Unification (\micro{}), \texttt{=°}] First-order syntactic unification with occurs check. Pairs,
  vectors and records are decomposed slot by slot, so a vector and a record with the same slots unify. When two
  variables meet, the one with the lower index is bound to the other. $\rightarrow$ \cref{ch:micro}.
\item[Uniform proof] A cut-free sequent proof in which every sequent with a non-atomic goal is the conclusion of the
  right-introduction rule for that goal's top-level connective. It formalises goal-directed
  search~\cite{MillerNadathurPfenningScedrov1991}. $\rightarrow$ \cref{ch:horn}; \book{73}.
\item[Universe index] (Nadathur, Jayaraman, Kwon) A numeric tag on constants and logic variables that records the
  $\forall$-nesting at their creation. Entering a universal goal increments the index. A variable may be bound only
  to terms whose constants have a tag no larger than its own, and the binding lowers the tags of the variables in
  the term~\cite{NadathurJayaramanKwon1995}. The book mentions the idea as numerical indices on quantified
  variables \book{225} and as universe level tags \book{227}. $\rightarrow$ \cref{ch:hounif,ch:lkanren}.
\item[Universe level] See \emph{Universe index} and \emph{Level of a unification variable}.
\end{glossarylist}

\subsection*{V}
\begin{glossarylist}
\item[Vacuous quantifier] A quantifier whose variable does not occur in its body. $\rightarrow$
  \cref{ch:lambdatree}; \book{195}.
\item[Value constructor] A constant of functional type used to build data, such as \lp{::}.
  $\rightarrow$ \cref{ch:intro}; \book{14}.
\item[Variable elimination (patterns)] The step of pattern unification that combines imitation, projection, the
  occurs check and pruning into one binding $F:=\lambda\bar c.\,t'$. $\rightarrow$ \cref{ch:hounif}; \book{223}.
\end{glossarylist}

\subsection*{W}
\begin{glossarylist}
\item[Walk] Resolving a variable through the substitution. \texttt{μkanren-state-find} is the shallow walk: it
  follows the chain of variable bindings until it reaches a non-variable or an unbound variable.
  \texttt{μkanren-state-find/value} is the deep walk: it also walks inside pairs, vectors and records. $\rightarrow$ \cref{ch:micro}.
\item[WAM] The Warren Abstract Machine, the classical implementation model of Prolog. $\rightarrow$ \book{74}.
\item[Working variable, reified variable] In \micro{}, a working variable has index $\ge 0$, lives during search
  and prints as \texttt{\_i}. A reified variable has a negative index, appears only in answers and prints as α, β,
  \dots{} $\rightarrow$ \cref{ch:micro}.
\end{glossarylist}

\endgroup

\section{A vocabulary bridge}\label{sec:glossary-bridge}

\Cref{tab:glossary-bridge} lines up the names that the three communities use for the same concept. Read it
across a row: a cell left empty means the concept has no counterpart in that column today.

\begin{table}[ht]
\centering\small
\caption{The same concepts under three names, and in the proposed module. The empty cells of the third
column are what \texttt{(aux kanren lambda)} supplies (\cref{ch:module}).}
\label{tab:glossary-bridge}
\begin{tabularx}{\textwidth}{@{}>{\raggedright}X>{\raggedright}X>{\raggedright}X>{\raggedright\arraybackslash}X@{}}
\toprule
Book (λProlog) & miniKanren literature & \micro{} today & \texttt{(aux kanren lambda)} \\
\midrule
logic variable, $\exists$ in goals (\texttt{sigma}) & \texttt{fresh} & \texttt{fresh°} & \texttt{fresh°}, with a level \\
backchaining on a clause & calling a relation & a \texttt{define-relation} call & also a \texttt{define-dynamic-relation} call \\
$\vee$ in goals (\texttt{;}) & \texttt{conde} & \texttt{or°}, \texttt{cond°} & unchanged \\
first-order unification with occurs check & \texttt{==} & \texttt{=°} & \texttt{=°}, through the \texttt{unify} and \texttt{post} hooks \\
\texttt{not}, \texttt{if} (via \texttt{!}) & \texttt{conda}, \texttt{condu} & \texttt{if°}, \texttt{take°} & unchanged \\
\rulename{augment}, \texttt{D => G} & \texttt{assume} (λKanren) & --- & \texttt{⊃°}, \texttt{define-dynamic-relation} \\
\rulename{generic}, \texttt{pi x\textbackslash{} G}, eigenvariable & \texttt{all} (λKanren), \texttt{eigen} & --- & \texttt{∀°}, with universe levels \\
$\lambda$-terms, \texttt{x\textbackslash{} t} & \texttt{tie}/\texttt{app} (λKanren), \texttt{tie} on noms (αKanren) & first-order encodings only (e.g.\ records with de Bruijn indices) & \texttt{ƛ}, \texttt{@} \\
pattern unification, raising, pruning & λKanren 2020 unifier; static closures (2021) & --- & the \texttt{unify} and \texttt{post} hooks with levels; non-patterns delayed in E \\
answer substitution & reified answer & \texttt{μkanren-project}, \texttt{°->list} & \texttt{°->list/datum}, \texttt{°->pretty} \\
\bottomrule
\end{tabularx}
\end{table}

\begin{kanren}[what the glossary says about the missing pieces]
The empty cells of \cref{tab:glossary-bridge} are the λProlog features that \micro{} lacks: a dynamic clause
database for AUGMENT, eigenvariables with a scope check for GENERIC, and terms with binders together with a
unifier that respects them. The entries \emph{Universe index}, \emph{Level of a unification variable}, \emph{Scope
check}, \emph{Higher-order pattern unification} and \emph{Delayed equations} list the known ways to provide them.
\Cref{ch:design} chooses among them and \cref{ch:module} implements the choice.
\end{kanren}

\begin{kanren}[three entries checked on today's engine]
The following file exercises the entries \emph{Interleaving}, \emph{Negation-as-failure} and \emph{Reification}.
The relation \texttt{num°} enumerates $0,1,2,\dots$.
\begin{scheme}
(import (aux kanren micro) (aux base) (aux stream))

; num° enumerates 0, 1, 2, ... (project° escapes to Scheme to add 1)
(define-relation (num° n)
  (or° (=° n 0) (fresh° (m) (num° m) (project° ((m* m)) (=° n (+ 1 m*))))))

; Interleaving: three infinite branches; the right-nested or° favours the first
(print (°->list/ground
         (take° 12 (fresh° (q n)
                     (or° (and° (num° n) (=° q (list 'a n)))
                          (and° (num° n) (=° q (list 'b n)))
                          (and° (num° n) (=° q (list 'c n))))))))

; Negation as failure: (take° 1 g) makes the test finite
(print (°->list/ground (fresh° (q) (=° q 1) (if° (take° 1 (num° q)) ✗° ✓°))))
(print (°->list/ground (fresh° (q) (=° q 1) (if° (=° q 2) ✗° ✓°))))
(print (°->list/ground (fresh° (q) (if° (=° q 2) ✗° ✓°))))

; Reification: free variables become α, β, ...; constraints first
(print (°->list #f (fresh° (q x) (=° q (list x 'a)) (≠° x 'b))))
\end{scheme}
\begin{answers}
((a 0) (b 0) (a 1) (c 0) (a 2) (b 1) (c 1) (a 3) (a 4) (b 2) (c 2) (a 5))
()
(1)
()
((λ (α) (begin (deny (equal? α (quote b)))) (cons α (cons (quote a) (quote ())))))
\end{answers}
Read the lines in order. (1) Of the first twelve answers, six come from branch \texttt{a} and three each from
\texttt{b} and \texttt{c}: the shares are $\frac12,\frac14,\frac14$, as the entry \emph{Interleaving} says.
(2) $q=1$ is a number, so \texttt{not (num° 1)} fails. If you drop the \texttt{take° 1}, the run never returns:
\texttt{num° 1} has only one answer, but the search for a second one never ends. \texttt{if°} continues every
answer of the test with \texttt{✗°}, and over a stream that never ends this continuation never finishes either. (3) and (4) show negation-as-failure on a ground goal and on a
non-ground one. With $q$ unbound, the test \texttt{(=° q 2)} succeeds by binding $q$, so the negation fails,
although $q=1$ would satisfy it. This is the classic unsoundness. (5) The reified answer is a closed
\texttt{λ}-expression whose body lists the constraint before the term.
\end{kanren}

\section{Summary}\label{sec:glossary-summary}
\begin{itemize}
\item Proof-theoretic terms (sequent, uniform proof, AUGMENT, GENERIC, eigenvariable) describe \emph{what} search
  must do. Unification terms (pattern, raising, pruning, level) describe \emph{how} binders and quantifiers are
  respected. miniKanren terms (goal, stream, interleaving, reification) describe the engine the ideas must fit.
\item Watch out for overloaded words: suspension, level, scope extrusion, fresh, cut, substitution.
\item Every entry points to the chapter that develops it and to the book page where it is defined. Use the
  pointers to return to the full treatment.
\end{itemize}
